%%
%% This is file `sample-acmsmall-conf.tex',
%% generated with the docstrip utility.
%%
%% The original source files were:
%%
%% samples.dtx  (with options: `all,proceedings,bibtex,acmsmall-conf')
%% 
%% IMPORTANT NOTICE:
%% 
%% For the copyright see the source file.
%% 
%% Any modified versions of this file must be renamed
%% with new filenames distinct from sample-acmsmall-conf.tex.
%% 
%% For distribution of the original source see the terms
%% for copying and modification in the file samples.dtx.
%% 
%% This generated file may be distributed as long as the
%% original source files, as listed above, are part of the
%% same distribution. (The sources need not necessarily be
%% in the same archive or directory.)
%%
%%
%% Commands for TeXCount
%TC:macro \cite [option:text,text]
%TC:macro \citep [option:text,text]
%TC:macro \citet [option:text,text]
%TC:envir table 0 1
%TC:envir table* 0 1
%TC:envir tabular [ignore] word
%TC:envir displaymath 0 word
%TC:envir math 0 word
%TC:envir comment 0 0
%%
%% The first command in your LaTeX source must be the \documentclass
%% command.
%%
%% For submission and review of your manuscript please change the
%% command to \documentclass[manuscript, screen, review]{acmart}.
%%
%% When submitting camera ready or to TAPS, please change the command
%% to \documentclass[sigconf]{acmart} or whichever template is required
%% for your publication.
%%
%%
\documentclass[acmsmall]{acmart}
%%
%% \BibTeX command to typeset BibTeX logo in the docs
\AtBeginDocument{%
  \providecommand\BibTeX{{%
    Bib\TeX}}}

%% Rights management information.  This information is sent to you
%% when you complete the rights form.  These commands have SAMPLE
%% values in them; it is your responsibility as an author to replace
%% the commands and values with those provided to you when you
%% complete the rights form.
\setcopyright{acmlicensed}
\copyrightyear{2018}
\acmYear{2018}
\acmDOI{XXXXXXX.XXXXXXX}
%% These commands are for a PROCEEDINGS abstract or paper.
\acmConference[Conference acronym 'XX]{Make sure to enter the correct
  conference title from your rights confirmation email}{June 03--05,
  2018}{Woodstock, NY}
%%
%%  Uncomment \acmBooktitle if the title of the proceedings is different
%%  from ``Proceedings of ...''!
%%
%%\acmBooktitle{Woodstock '18: ACM Symposium on Neural Gaze Detection,
%%  June 03--05, 2018, Woodstock, NY}
\acmISBN{978-1-4503-XXXX-X/2018/06}


%%
%% Submission ID.
%% Use this when submitting an article to a sponsored event. You'll
%% receive a unique submission ID from the organizers
%% of the event, and this ID should be used as the parameter to this command.
%%\acmSubmissionID{123-A56-BU3}

%%
%% For managing citations, it is recommended to use bibliography
%% files in BibTeX format.
%%
%% You can then either use BibTeX with the ACM-Reference-Format style,
%% or BibLaTeX with the acmnumeric or acmauthoryear sytles, that include
%% support for advanced citation of software artefact from the
%% biblatex-software package, also separately available on CTAN.
%%
%% Look at the sample-*-biblatex.tex files for templates showcasing
%% the biblatex styles.
%%

%%
%% The majority of ACM publications use numbered citations and
%% references.  The command \citestyle{authoryear} switches to the
%% "author year" style.
%%
%% If you are preparing content for an event
%% sponsored by ACM SIGGRAPH, you must use the "author year" style of
%% citations and references.
%% Uncommenting
%% the next command will enable that style.
%%\citestyle{acmauthoryear}


\usepackage{booktabs}
\usepackage{caption}
\usepackage{threeparttable}
\captionsetup[table]{skip=5pt} % Ajusta el 'skip' a los puntos que desees
%%
%% end of the preamble, start of the body of the document source.
\begin{document}

%%
%% The "title" command has an optional parameter,
%% allowing the author to define a "short title" to be used in page headers.
\title{Monitoring and Understanding Notebook-Based Learning with Jupyter - Appendix}
\subtitle{A Dashboard for Student Self-Assessment - Appendix}

%%
%% The "author" command and its associated commands are used to define
%% the authors and their affiliations.
%% Of note is the shared affiliation of the first two authors, and the
%% "authornote" and "authornotemark" commands
%% used to denote shared contribution to the research.
\author{Arturo Olivares Martos}
\email{arturoolivares@correo.ugr.es}
\email{arturo.olivares@stud.uni-due.de}
\affiliation{%
  \institution{Universidad de Granada}
  \city{Granada}
  \country{Spain}
}
\affiliation{%
  \institution{University of Duisburg-Essen}
  \city{Duisburg}
  \country{Germany}
}

\author{Lisa Pawlowski}
\email{lisa.pawlowski.99@stud.uni-due.de}
\affiliation{%
  \institution{University of Duisburg-Essen}
  \city{Duisburg}
  \country{Germany}
}

\author{Mohamed Abdelmagied}
\email{mohamed.abdelmagied@stud.uni-due.de}
\affiliation{%
  \institution{University of Duisburg-Essen}
  \city{Duisburg}
  \country{Germany}
}

\author{Kaimao Sheng}
\email{kaimao.sheng@uni-due.de}
\affiliation{%
  \institution{University of Duisburg-Essen}
  \city{Duisburg}
  \country{Germany}
}

\author{Irene-Angelica Chounta}
\email{irene-angelica.chounta@uni-due.de}
\affiliation{%
  \institution{University of Duisburg-Essen}
  \city{Duisburg}
  \country{Germany}
}

%%
%% By default, the full list of authors will be used in the page
%% headers. Often, this list is too long, and will overlap
%% other information printed in the page headers. This command allows
%% the author to define a more concise list
%% of authors' names for this purpose.
%\renewcommand{\shortauthors}{Trovato et al.}

%%
%% The abstract is a short summary of the work to be presented in the
%% article.
%%
%% This command processes the author and affiliation and title
%% information and builds the first part of the formatted document.
\maketitle
%% MUC OUTLINE
%% Lisa and Arturo, please fill out the following points

\appendix
\counterwithin{table}{section}
\counterwithin{figure}{section}

\section{Questionnaire Items}
\label{sec:appendix_questionnaires}
This appendix provides the exact wording of the items used in the evaluation survey to measure system usability, technology acceptance, and perceived cognitive workload.

\subsection{System Usability Scale (SUS)}
Participants rated the 10 items of Table~\ref{tab:app_sus} on a 5-point Likert scale (1 = Strongly Disagree to 5 = Strongly Agree). 

\begin{table}[htpb]
  \centering
  \caption{Items of the System Usability Scale (SUS)}
  \label{tab:app_sus}
  \begin{tabular}{lp{0.85\linewidth}}
    \toprule
    \textbf{Item} & \textbf{Statement} \\
    \midrule
    SUS01 & I think I would like to use this Dashboard frequently. \\
    SUS02 & I found this Jupyter Dashboard unnecessarily complex. \\
    SUS03 & I thought the Jupyter Dashboard was easy to use. \\
    SUS04 & I think that I would need the support of a technical person to use the Jupyter Dashboard. \\
    SUS05 & I found the various functions in this dashboard were well integrated. \\
    SUS06 & I thought there was too much inconsistency in this dashboard. \\
    SUS07 & I would imagine that most people would learn to use this dashboard very quickly. \\
    SUS08 & I found this dashboard very awkward to use. \\
    SUS09 & I felt very confident using this Jupyter Dashboard. \\
    SUS10 & I needed to learn a lot of things before I could get going with this dashboard. \\
    \bottomrule
  \end{tabular}
\end{table}

\subsection{Technology Acceptance Model (TAM)}
Participants rated the items of Table~\ref{tab:app_tam} on a 7-point Likert scale (Extremely Unlikely to Extremely Likely). The construct was divided into Perceived Usefulness (PU) and Perceived Ease of Use (PEOU).

\begin{table}[htpb]
  \centering
  \caption{Items of the Technology Acceptance Model (TAM)}
  \label{tab:app_tam}
  \begin{tabular}{lp{0.85\linewidth}}
    \toprule
    \textbf{Construct} & \textbf{Statement} \\
    \midrule
    \textbf{PU} & \textit{Using the Jupyter Dashboard to assist my learning progress would...} \\
    PU01 & ...enable me to accomplish tasks more frequently. \\
    PU02 & ...improve my learning performance. \\
    PU03 & ...increase my productivity. \\
    PU04 & ...enhance my effectiveness during the task. \\
    PU05 & ...make it easier to do my tasks. \\
    PU06 & ...be useful. \\
    \midrule
    \textbf{PEOU} & \textit{How much do you agree with each statement?} \\
    PEOU01 & Learning to operate this Dashboard would be easy for me. \\
    PEOU02 & I would find it easy to get this Dashboard to do what I want it to do. \\
    PEOU03 & My interaction with the Dashboard would be clear and understandable. \\
    PEOU04 & I would find this Dashboard would be clear and understandable. \\
    PEOU05 & It would be easy for me to become skillful at using this Dashboard. \\
    PEOU06 & I would find this Dashboard easy to use. \\
    \bottomrule
  \end{tabular}
\end{table}

\subsection{NASA Task Load Index (NASA-TLX)}
Participants rated the items of Table~\ref{tab:app_nasatlx} on a continuous scale from 0 to 20 to measure their perceived workload during the task.

\begin{table}[htpb]
  \centering
  \caption{Items of the NASA-TLX}
  \label{tab:app_nasatlx}
  \begin{tabular}{lp{0.85\linewidth}}
    \toprule
    \textbf{Dimension} & \textbf{Question} \\
    \midrule
    Mental Demand & How mentally demanding was the task? \\
    Physical Demand & How physically demanding was the task? \\
    Temporal Demand & How hurried or rushed was the pace of the task? \\
    Performance & How successful were you in accomplishing what you were asked to do? \\
    Effort & How hard did you have to work to accomplish your level of performance? \\
    Frustration & How insecure, discouraged, irritated, stressed, and annoyed were you? \\
    \bottomrule
  \end{tabular}
\end{table}



\end{document}
\endinput
%%
%% End of file `sample-acmsmall-conf.tex'.
